\documentclass[10pt,a4paper]{amsart}
\usepackage[margin=19mm]{geometry}
\usepackage{amsmath,amssymb,mathtools}
\usepackage[scr=rsfs,cal=euler]{mathalfa}
\usepackage{xfrac}
\usepackage[unicode,hidelinks,bookmarks=false]{hyperref}
\newcommand{\ie}{i.e.\ }
\newcommand{\cf}{cf.\ }
\newcommand{\resp}{respectively\ }
\newcommand{\Subsection}{\subsection}
\providecommand{\nobreakdash}{\nobreak\hbox{-}}
\setlength{\emergencystretch}{2em}
\multlinegap0pt
\usepackage{bm}
\usepackage{bbm}
\usepackage{dsfont}
\usepackage{xspace}
\usepackage{cancel}
\usepackage{mathtools}
\usepackage[shortlabels]{enumitem}
\usepackage{amsthm}
\makeatletter
\@ifundefined{theorem}{\newtheorem{theorem}{Theorem}}{}
\@ifundefined{lemma}{\newtheorem{lemma}[theorem]{Lemma}}{}
\@ifundefined{proposition}{\newtheorem{proposition}[theorem]{Proposition}}{}
\makeatother
\title[Characterisations of Kullback--Leibler approximation by fini]{Characterisations of Kullback--Leibler approximation by finite Gaussian mixtures}
\author{Hien Duy Nguyen}
\date{Source: arXiv:2604.10899v1 (CC BY 4.0)}
\begin{document}
\maketitle

\noindent\emph{Source and licence.} Characterisations of Kullback--Leibler approximation by finite Gaussian mixtures. Version arXiv:2604.10899v1 (CC BY 4.0), distributed under the Creative Commons Attribution 4.0 International licence, as recorded in the arXiv licence metadata for this exact version and shown on the abstract page https://arxiv.org/abs/2604.10899v1. Full rights notice: licenses/arxiv-2604.10899-CC-BY-4.0.txt.\par\smallskip

\noindent\textit{Editorial scope.} Counted unit: the Theorem whose source label is \texttt{thm:entclass-UI} in the article (source0.tex lines 311--319, source label \texttt{thm:entclass-UI}), delivered together with the article's own complete original proof of it (source0.tex lines 321--491, one \texttt{proof} environment of 171 lines). No mathematics is written, repaired or re-derived: statement and proof are the article's own text line for line. Required context retained uncounted: prop:UIKL (lines 136--152); thm:necessity (lines 238--249); lem:compactapprox (lines 1349--1355). Every definition, display and label of those lines is retained unchanged. Omitted from this package and not redistributed: 1--135, 153--237, 250--310, 320--320, 492--1348, 1356--1496. Nothing inside a retained excerpt is omitted or altered: every retained body is delivered exactly as the article's lines are written, so no omission receipt of any kind accompanies this package. Citations: the retained excerpts carry no citation command at all, so no bibliography entry of the article is transcribed and no reference is redistributed. Reference adaptation: none was needed and none was made; every reference inside the delivered unit resolves inside it, so no unresolved reference or citation is produced. Printed slips kept literal: no printed word, symbol, hypothesis or number of the counted statement or of its proof was altered. Layout and class: the article's own class is \texttt{article} with packages including geometry, amsmath, amsthm, amssymb, mathtools, bm, fontenc, lmodern, microtype, natbib, hyperref, enumitem, mathrsfs; the wrapper is the lane's pinned \texttt{amsart} editorial wrapper at 10pt on A4, which restates the theorem-like environments the retained text needs and adds the article's own macro definitions that the retained text needs (none), with extra wrapper packages (bm, bbm, dsfont, xspace, cancel, mathtools, enumitem). The delivered numbering is the unit's own. Assets: no retained span contains a figure environment, a graphics-inclusion command, a PDF-page inclusion, a table or a TikZ picture; no image file is copied into this package, the package declares no asset dependency and no image file is redistributed with it. Licence: version arXiv:2604.10899v1 of this article is distributed under the Creative Commons Attribution 4.0 International licence, as recorded in the arXiv licence metadata for this exact version and shown on the abstract page https://arxiv.org/abs/2604.10899v1. A full rights notice is delivered at licenses/arxiv-2604.10899-CC-BY-4.0.txt. Characters: every retained excerpt is pure ASCII, so the delivered engine, XeLaTeX with its default Latin Modern text font, reports no missing character.
\begin{proposition}
\label{prop:UIKL}
Let $(g_m)_{m\ge1}$ be densities on $\mathbb{R}^d$ such that:
\begin{enumerate}[label=\textnormal{(\roman*)}]
    \item $g_m(x)>0$ for $F$-almost every $x$;
    \item $f(x)/g_m(x)\to1$ for $F$-almost every $x$;
    \item the family $\{L_m\}_{m\ge1}$, where
    \[
    L_m=\log_+\!\left(\frac{f}{g_m}\right),
    \]
    is uniformly integrable under $F$.
\end{enumerate}
Then
\[
\mathrm{KL}(f\|g_m)\to0.
\]
\end{proposition}

\begin{theorem}[Universal necessity of the second moment]
\label{thm:necessity}
Let $f$ be any density on $\mathbb{R}^d$. If there exists $g\in\mathcal M$ such that $\mathrm{KL}(f\|g)<\infty$, then
\[
\int_{\mathbb{R}^d}\left\lVert x \right\rVert^2 f(x)\,\mathrm{d}x<\infty.
\]
In particular, if
\[
\inf_{g\in\mathcal M}\mathrm{KL}(f\|g)=0,
\]
then the same second-moment conclusion holds.
\end{theorem}

\begin{lemma}[Uniform approximation on a fixed ball]
\label{lem:compactapprox}
Let $f$ be a continuous probability density on $\mathbb{R}^d$. For every $R>0$ and every $\eta>0$, there exists a finite GMM $h\in\mathcal M$ such that
\[
\sup_{x\in B_R} |h(x)-f(x)|<\eta.
\]
\end{lemma}

\begin{theorem}[Equivalence on the finite log-moment class]
\label{thm:entclass-UI}
For every $f\in\mathcal F_{\mathrm{Ent}}$,
\[
\inf_{g\in\mathcal M}\mathrm{KL}(f\|g)=0
\quad\Longleftrightarrow\quad
\int_{\mathbb{R}^d}\left\lVert x \right\rVert^2f(x)\,\mathrm{d}x<\infty.
\]
\end{theorem}

\begin{proof}
Assume first that
\[
\int_{\mathbb{R}^d}\left\lVert x \right\rVert^2f(x)\,\mathrm{d}x<\infty.
\]
For $R>0$, define the tail quantities
\[
\alpha_R=\int_{B_{R}^c} f(x)\,\mathrm{d}x,
\qquad
\beta_R=\int_{B_{R}^c} f(x)\log_+ f(x)\,\mathrm{d}x,
\qquad
\mu_{2,R}=\int_{B_{R}^c}\left\lVert x \right\rVert^2 f(x)\,\mathrm{d}x.
\]
Since $f\log_+f\in \mathcal L^1(\mathbb{R}^d)$ and $\left\lVert \cdot \right\rVert^2 f\in \mathcal L^1(\mathbb{R}^d)$, we have
\[
\alpha_R\to0,
\qquad
\beta_R\to0,
\qquad
\mu_{2,R}\to0
\qquad (R\to\infty).
\]
Moreover, because $\alpha_R\to0$,
\[
\alpha_R\log\!\left(\frac1{\alpha_R}\right)\to0.
\]
Also $\alpha_R\log R\to0$: if $s_2=\int \left\lVert x \right\rVert^2f(x)\,\mathrm{d}x<\infty$, then Markov's inequality gives $\alpha_R\le s_2/R^2$, hence
\[
0\le \alpha_R\log R\le s_2\frac{\log R}{R^2}\to0.
\]

Choose a sequence $(R_m)_{m\ge1}$ such that $R_m\uparrow\infty$, $R_m\ge1$ for every $m$, and, for every $m\ge1$,
\begin{equation}
\label{eq:Rmchoice}
\begin{aligned}
\alpha_{R_m}&<\tfrac14,
&\beta_{R_m}&<2^{-m-4},
&\alpha_{R_m}\log\!\left(\frac1{\alpha_{R_m}}\right)&<2^{-m-4},\\
\alpha_{R_m}\frac d2\log(2\pi R_m^2)&<2^{-m-4},
&\frac{\mu_{2,R_m}}{2R_m^2}&<2^{-m-4}.
\end{aligned}
\end{equation}
This is possible by the convergences above.

For each $m$, continuity and strict positivity imply
\[
m_m=\inf_{x\in B_{R_m}} f(x)>0.
\]
Apply Lemma~\ref{lem:compactapprox} with radius $R_m$ and tolerance
\[
\eta_m=2^{-m-3}m_m
\]
to obtain a finite GMM $h_m\in\mathcal M$ such that
\begin{equation}
\label{eq:hmclose}
\sup_{x\in B_{R_m}}|h_m(x)-f(x)|<\eta_m.
\end{equation}
In particular,
\begin{equation}
\label{eq:hmlower}
h_m(x)\ge f(x)-\eta_m\ge m_m-2^{-m-3}m_m\ge \frac{m_m}{2}
\qquad \text{for all }x\in B_{R_m}.
\end{equation}
Now define
\[
\phi_m(x)=\varphi(x;0,R_m^2 I_d),
\qquad
g_m(x)=(1-\alpha_{R_m})h_m(x)+\alpha_{R_m}\phi_m(x).
\]
Then each $g_m$ is a finite GMM, hence $g_m\in\mathcal M$.

Fix $x\in\mathbb{R}^d$. For all sufficiently large $m$, $x\in B_{R_m}$, so by \eqref{eq:hmclose},
\[
|h_m(x)-f(x)|\le \eta_m=2^{-m-3}m_m\le 2^{-m-3}f(x)\to0.
\]
Also $\alpha_{R_m}\to0$, and
\[
\phi_m(x)=(2\pi R_m^2)^{-d/2}\exp\!\left(-\frac{\left\lVert x \right\rVert^2}{2R_m^2}\right)\to0.
\]
Therefore
\[
|g_m(x)-f(x)|
\le (1-\alpha_{R_m})|h_m(x)-f(x)| + \alpha_{R_m}f(x)+\alpha_{R_m}\phi_m(x)\to0.
\]
Since $f(x)>0$ for every $x$, it follows that $f(x)/g_m(x)\to1$ pointwise on $\mathbb{R}^d$.

We next estimate $L_m=\log_+(f/g_m)$ in $\mathcal L^1(F)$. Since $g_m\ge (1-\alpha_{R_m})h_m$, we have for $x\in B_{R_m}$,
\[
\log\!\left(\frac{f(x)}{g_m(x)}\right)
\le \log\!\left(\frac{f(x)}{(1-\alpha_{R_m})h_m(x)}\right)
=\log\!\left(\frac{f(x)}{h_m(x)}\right)+\log\!\left(\frac1{1-\alpha_{R_m}}\right).
\]
Hence
\[
L_m(x)\le \left|\log\!\left(\frac{f(x)}{h_m(x)}\right)\right|+\log\!\left(\frac1{1-\alpha_{R_m}}\right).
\]
On $B_{R_m}$, both $f$ and $h_m$ are at least $m_m/2$ by \eqref{eq:hmlower}; therefore
\[
\left|\log\!\left(\frac{f(x)}{h_m(x)}\right)\right|
\le \frac{|f(x)-h_m(x)|}{m_m/2}
\le \frac{2\eta_m}{m_m}=2^{-m-2}.
\]
Using $\log(1/(1-u))\le 2u$ for $u\in[0,1/2]$ and \eqref{eq:Rmchoice}, we conclude that on $B_{R_m}$,
\begin{equation}
\label{eq:compactL1}
\int_{B_{R_m}}L_m\,\mathrm{d}F \le 2^{-m-2}+2\alpha_{R_m}.
\end{equation}

Since $R_m\ge1$ and $x\in B_{R_m}^c$, we have
\[
\phi_m(x)
=(2\pi R_m^2)^{-d/2}\exp\!\left(-\frac{\left\lVert x \right\rVert^2}{2R_m^2}\right)
\le (2\pi R_m^2)^{-d/2}e^{-1/2}
\le (2\pi)^{-d/2}e^{-1/2}<1.
\]
Because $\alpha_{R_m}<1/4$, it follows that
\[
\log\!\left(\frac{1}{\alpha_{R_m}\phi_m(x)}\right)\ge0.
\]
Since $g_m\ge \alpha_{R_m}\phi_m$ on all of $\mathbb{R}^d$, we therefore obtain on $B_{R_m}^c$,
\[
L_m(x)
\le \left(\log f(x)+\log\!\left(\frac1{\alpha_{R_m}\phi_m(x)}\right)\right)_+
\le \log_+ f(x)+\log\!\left(\frac1{\alpha_{R_m}}\right)-\log\phi_m(x).
\]
Because
\[
-\log\phi_m(x)=\frac d2\log(2\pi R_m^2)+\frac{\left\lVert x \right\rVert^2}{2R_m^2},
\]
integration over $B_{R_m}^c$ yields
\begin{equation}
\label{eq:tailL1}
\int_{B_{R_m}^c}L_m\,\mathrm{d}F
\le \beta_{R_m}+\alpha_{R_m}\log\!\left(\frac1{\alpha_{R_m}}\right)+\alpha_{R_m}\frac d2\log(2\pi R_m^2)+\frac{\mu_{2,R_m}}{2R_m^2}.
\end{equation}
By \eqref{eq:Rmchoice}, the right-hand side is at most $2^{-m-2}$.

Combining \eqref{eq:compactL1} and \eqref{eq:tailL1}, and using again \eqref{eq:Rmchoice}, we obtain
\[
\int L_m\,\mathrm{d}F \le 2^{-m-2}+2\alpha_{R_m}+2^{-m-2}\to0.
\]
In particular, $L_m\to0$ in $\mathcal L^1(F)$. To verify uniform integrability directly, let $\varepsilon>0$. Choose $m_0$ so large that
\[
\int L_m\,\mathrm{d}F<\varepsilon
\qquad\text{for all }m\ge m_0.
\]
For the finitely many indices $1\le m<m_0$, choose $M<\infty$ such that
\[
\int_{\{L_m>M\}}L_m\,\mathrm{d}F<\varepsilon
\qquad\text{for all }1\le m<m_0.
\]
Then, for $m\ge m_0$,
\[
\int_{\{L_m>M\}}L_m\,\mathrm{d}F\le \int L_m\,\mathrm{d}F<\varepsilon,
\]
and for $1\le m<m_0$ the same bound holds by construction. Hence
\[
\sup_{m\ge1}\int_{\{L_m>M\}}L_m\,\mathrm{d}F<\varepsilon,
\]
so the family $\{L_m\}_{m\ge1}$ is uniformly integrable. Proposition~\ref{prop:UIKL} therefore gives
\[
\mathrm{KL}(f\|g_m)\to0.
\]
This proves that finite second-moment existence implies KL-approximability and yields the approximating sequence. The implication
\[
\inf_{g\in\mathcal M}\mathrm{KL}(f\|g)=0
\quad\Longrightarrow\quad
\int_{\mathbb{R}^d}\left\lVert x \right\rVert^2f(x)\,\mathrm{d}x<\infty
\]
follows from Theorem~\ref{thm:necessity}.
\end{proof}

\end{document}
